\documentclass[10pt,a4paper]{article}
\usepackage[margin=2cm]{geometry}
\usepackage{multicol}
\usepackage{times}
\usepackage{graphicx}
\usepackage{float}
\usepackage{amsmath}
\usepackage{booktabs}
\usepackage{caption}

\title{\textbf{Head Circumference Estimation from Ultrasound Images Using Convolutional Neural Networks}}
\author{Dinh Thanh Hai -- 22BA13119}
\date{}

\begin{document}
\maketitle

\begin{multicols}{2}

\section*{Abstract}
Ultrasound imaging is widely used to monitor fetal development during pregnancy.
One important measurement is the head circumference (HC), which helps doctors
evaluate the growth of the fetus.
In this study, a convolutional neural network (CNN) is used to estimate the head
circumference directly from ultrasound images.
The model is trained and evaluated using Mean Absolute Error (MAE).
The experimental results show that the proposed method can predict head
circumference with good accuracy.

\section*{I. Introduction}
Medical imaging plays an important role in prenatal healthcare.
Ultrasound images are commonly used because they are safe and non-invasive.
Head circumference is a key indicator of fetal growth and brain development.
Traditionally, this measurement is performed manually by medical experts.

Recently, deep learning methods have shown strong performance in medical image
analysis.
Therefore, this study aims to use a convolutional neural network to automatically
predict the head circumference from ultrasound images.
This approach can help reduce manual effort and improve measurement consistency.

\section*{II. Dataset Description}
The dataset used in this study consists of ultrasound images and corresponding
head circumference values.
The data is divided into a training set and a validation set.

\subsection*{A. Data Structure}
Each data sample includes:
\begin{itemize}
    \item An ultrasound image file.
    \item Pixel size information in millimeters.
    \item The ground truth head circumference value in millimeters.
\end{itemize}

A total of 999 images are provided, where 799 images are used for training and
200 images are used for validation.

\subsection*{B. Summary Statistics}
The head circumference values range from approximately 44 mm to 346 mm.
Most samples are concentrated around the middle range, which represents normal
fetal development stages.

\subsection*{C. Data Distribution}
Figure~1 shows the distribution of head circumference values in the dataset.
The histogram indicates that the data is not uniformly distributed.
However, it still covers a wide range of values, which is useful for training
a regression model.

\begin{figure}[H]
    \centering
    \includegraphics[width=\linewidth]{download-7.png}
    \caption{Distribution of head circumference values}
\end{figure}

\section*{III. Implementation}
The model is implemented using Python and the PyTorch deep learning framework.
Google Colab is used for training with GPU acceleration.

\subsection*{A. Preprocessing}
All images are resized and normalized before being fed into the network.
The dataset is loaded using a custom PyTorch Dataset class.
The training set contains 799 images, while the validation set contains 200 images.

\subsection*{B. Model Selection}
A convolutional neural network based on a pretrained ResNet architecture is used.
The pretrained weights help the model learn faster and improve performance.
The final fully connected layer is modified to output a single continuous value
for head circumference prediction.

\section*{IV. Training Setup}
The model is trained using the Mean Absolute Error (MAE) loss function.
The Adam optimizer is applied to update the model parameters.
Training is performed for multiple epochs until the validation error stabilizes.

The dataset split is shown below:
\begin{itemize}
    \item Training set: 799 images
    \item Validation set: 200 images
\end{itemize}

\section*{V. Results and Discussion}
After training, the model is evaluated on the validation set.
The final Mean Absolute Error obtained is 10.97 mm.

Figure~2 shows the relationship between predicted values and ground truth values.
Most points are close to the diagonal line, which represents perfect prediction.
This indicates that the model is able to learn the relationship between ultrasound
images and head circumference.

\begin{figure}[H]
    \centering
    \includegraphics[width=\linewidth]{download-8.png}
    \caption{Predicted head circumference versus ground truth}
\end{figure}

Although the results are good, some prediction errors still exist.
This may be caused by image noise, variation in fetal position, or limited data.

\section*{VI. Conclusion}
This study demonstrates that a convolutional neural network can be used to
estimate fetal head circumference from ultrasound images.
The proposed model achieves a validation MAE of 10.97 mm, which shows good
prediction accuracy.

In the future, performance can be improved by using more training data,
stronger data augmentation, or more advanced neural network architectures.
This approach has the potential to support doctors in prenatal healthcare.

\end{multicols}

\end{document}
